% 14-2008-middle-east — 2008 Test Run, Middle East (dust): daily regional-mean SW AOD
% (four bands) and LW optical depth (four bands) on the top row, CAM7 vs
% AERONET daily AOD in the visible (546 vs 500 nm) and near-infrared
% (842 vs 870 nm) on the bottom row. Slides 13-24 are grouped by dominant
% aerosol type (dust, smoke, pollution, ocean outflow) and within each type
% by decreasing AERONET 500 nm annual mean (analyses/statistics/.../README.md).
% Sources: DAVINCI regional-aod and regional-aeronet analyses
% (analyses/davinci/), copied under graphics/2008-regional/.
% All four plots are 656x296pt with ~7pt margins: no trim; \regpanel
% (preamble.tex) sets each at 0.48 of the column width, ~2.7cm tall, so
% two rows fit under \bodyheight.
\begin{frame}[t]
\frametitle{2008 Middle East}
\centering
\regpanel{2008-regional/middle-east/middle_east_shortwave_aod.pdf}\hfill
\regpanel{2008-regional/middle-east/middle_east_longwave_optical_depth.pdf}\par
\vspace{0.15cm}
\regpanel{2008-regional/middle-east/middle_east_aeronet_visible.pdf}\hfill
\regpanel{2008-regional/middle-east/middle_east_aeronet_near_infrared.pdf}\par
\vspace{0.1cm}
{\tiny\setlength{\tabcolsep}{5pt}
\begin{tabular}{@{}lrrrrrrr@{}}
\toprule
CAM7 vs AERONET & Days & AERONET & CAM7 & Bias & Mean error & RMSE & R \\
\midrule
Visible (546 vs 500 nm) & 365 & 0.368 & 0.313 & $-0.055$ & 0.070 & 0.094 & 0.91 \\
Near-infrared (842 vs 870 nm) & 365 & 0.298 & 0.206 & $-0.093$ & 0.097 & 0.129 & 0.93 \\
\bottomrule
\end{tabular}\par}
\end{frame}
